\section{Motivating Examples}

\begin{figure}[!t]
    \begin{minted}[xleftmargin=5pt, numbersep=4pt, linenos = true, fontsize = \footnotesize, escapeinside=??]{OCaml}
    let rec merge (l1: int list) (l2: int list) : int list =
      match l1, l2 with
      | [], l2 -> l2
      | l1, [] -> l1
      | h1 :: t1, h2 :: t2 ->
        if h1 < h2 then h1 :: h2 :: merge t1 t2
        else if h1 > h2 then h2 :: h1 :: merge t1 t2
        else h1 :: h2 :: merge t1 t2
    
    let rec mergesort (split: int list -> int list * int list) (l: int list) =
        match l with
        | [] | [_] -> l
        | _ ->
            let l1, l2 = split l in
            merge (mergesort l1) (mergesort l2)
    \end{minted}
    \caption{\Code{mergesort} implements the merge sort algorithm: it splits the list, recursively sorts the sublists, and merges the results (lines $14$--$15$). \Code{mergesort} is parameterized by the input \Code{split} function, which can use different splitting strategies.
    The $\Code{merge}$ function is supposed to merge two sorted lists. If either of the input lists is empty, \Code{merge} returns the other (lines $3$--$4$). Otherwise, the function merges the lists recursively (lines $6$--$8$). However, this implementation is buggy: in the case $\Code{h1} < \Code{h2}$, it places $\Code{h2}$ before the result of merging $\Code{t1}$ and $\Code{t2}$, incorrectly assuming that $\Code{h2}$ is smaller than all elements in $\Code{t1}$. }
    \label{fig:motivation}
\end{figure}

\paragraph{Proving Incorrectness} We attempt to prove the incorrectness of \Code{mergesort}, as shown in Figure \ref{fig:motivation}; that is, it violates the sorting property.
\begin{align*}
    \text{sorting property: } & \forall l\;l'.\ l' = \Code{mergesort}\ l \impl \I{sorted}(l')
    \\ \text{incorrectness: } & \exists l\;l'.\ l' = \Code{mergesort}\ l \land \neg \I{sorted}(l')
\end{align*}\noindent
More precisely, we would like to know under what conditions (e.g., shape of input list, properties of the \Code{split} function) the \Code{mergesort} function is incorrect.

\paragraph{Example} For input $\Code{l1 = [1;2],\ l2 = [5;6]}$, the \Code{merge} function will produce the unsorted output $\Code{[1;5;2;6]}$. If we choose the second implementation of the \Code{split} function shown in Figure \ref{fig:split}, the \Code{mergesort} function will produce the unsorted list from the input $\Code{[2;1;6;5]}$.
\begin{figure}[!t]
    \begin{minted}[xleftmargin=5pt, numbersep=4pt, linenos = true, fontsize = \footnotesize, escapeinside=??]{OCaml}
    (* the first list will not be less than the second list *)
    let split (l: int list) : int list * int list =
        let rec aux l1 l2 = function
          | [] -> (l1, l2)
          | [x] -> (x :: l1, l2)
          | x :: y :: rest ->
              aux (x :: l1) (y :: l2) rest
        in
        aux [] [] l

    (* the first list will be no longer than the second list *)
    let split (l: int list) : int list * int list =
        let n = List.length l in
        let n1 = n / 2 in
        let n2 = n - n1 in
        (List.sublist l 0 n1, List.sublist l n1 n)

    (* random pick an element as pivot and split the list into two parts ... *)
    \end{minted}
    \caption{Two possible implementations of the \Code{split} function.}
    \label{fig:split}
\end{figure}

\paragraph{Challenges} There are several challenges in proving incorrectness.
\begin{enumerate}
    \item Higher-order functions, which prevent traditional testing and symbolic execution approaches. We also want to reject some trivial invalid input functions.
    \item Compositional reasoning: we would like to find a precise incorrectness precondition for \Code{mergesort} to enable further proofs over different implementations of \Code{split}. For example, the following specification is not allowed:
    \begin{align*}
        \Code{split} : \ort{\List[\Int]}{\top} \sarr \urt{\List[\Int]}{\top} \times \urt{\List[\Int]}{\top}
    \end{align*}\noindent
    where $\Code{split}$ should follow some basic properties, e.g., preserve membership of the input list.
    Then the incorrectness precondition of $\Code{merge}$ should also be precise; i.e., it cannot just be 
    \begin{align*}
        \Code{merge} : \urt{\List[\Int]}{\top} \sarr \urt{\List[\Int]}{\top} \sarr \Code{unsorted}
    \end{align*}\noindent
    \item Recursive functions, where the incorrectness execution path involves multiple recursive calls. For some calls, we need a correctness specification; for others, we need an incorrectness specification.
\end{enumerate}

\paragraph{Solution: Using Over- and Under- Approximation}
\begin{enumerate}
    \item Higher-order function: We provide an ``over-style'' specification of the input function:
    \begin{align*}
        &\Code{mergesort}:
        \\&\quad \ol{\Code{split}}{:}(\ol{l}{:}\rt{il}{\top} \sarr \ol{\Code{l1}}{:}\rt{il}{\top} \times \rt{il}{\I{len}(\vnu) = \I{len}(\Code{l1}) + \I{len}(\vnu) \land \I{mem}\ldots}) \sarr
        \\&\quad \ol{\Code{l}}{:}\rt{il}{\top} \sarr
        \\&\quad \rt{il}{\I{sorted}(\vnu)}
    \end{align*}\noindent
    Then we ``negate'' it to get the ``under-style'' specification:
    \begin{align*}
        &\Code{mergesort}:
        \\&\quad \ul{\Code{split}}{:}(\ol{l}{:}\rt{il}{\top} \sarr \ol{\Code{l1}}{:}\rt{il}{\top} \times \rt{il}{\I{len}(\vnu) = \I{len}(\Code{l1}) + \I{len}(\vnu) \land \I{mem}\ldots}) \sarr
        \\&\quad \ul{\Code{l}}{:}\rt{il}{\top} \sarr
        \\&\quad \rt{il}{\neg \I{sorted}(\vnu)}
    \end{align*}\noindent
    Note that the function type of $\Code{split}$ is not flipped (it is still an over-style specification), but the way we treat the parameter $\Code{split}$ is changed to under-style (similar to higher-order model checking). Our proof should show that the program can be typechecked against some supertype of this type. Similarly, the $\Code{merge}$ function should have an over-style specification:
    \begin{align*}
        &\Code{merge}: \ol{\Code{l1}}{:}\rt{il}{\I{sorted}(\vnu)} \sarr \ol{\Code{l2}}{:}\rt{il}{\I{sorted}(\vnu)} \sarr \rt{il}{\top}
    \end{align*}\noindent
    \item Mixed specification: The $\Code{merge}$ function requires two parameters; if both of them are underapproximated, it is too specific and hard to match the actual implementation. For example, the two implementations of the $\Code{split}$ function shown in Figure \ref{fig:split} have different guarantees on the lengths of the output lists. However, the following specification is more general:
    \begin{align*}
        &\Code{merge}:
        \\&\quad \ol{\Code{l1}}{:}\rt{il}{\I{len}(\vnu) \geq 2 \land \I{sorted}(\vnu)} \sarr \ul{\Code{l2}}{:}\rt{il}{\I{len}(\vnu) = 2 \land \I{sorted}(\vnu)} \sarr \rt{il}{\neg \I{sorted}(\vnu)}
        \\&\quad \ol{\Code{l1}}{:}\rt{il}{\I{len}(\vnu) \geq 2 \land \I{sorted}(\vnu)} \sarr \ul{\Code{l2}}{:}\unionty_{i{:}\Int}\rt{il}{\I{len}(\vnu) = i \land i \geq 2 \land \I{sorted}(\vnu)} \sarr \rt{il}{\neg \I{sorted}(\vnu)}
    \end{align*}\noindent
    This can lead to a more general specification of the $\Code{split}$ function (i.e., the incorrectness precondition):
    \begin{align*}
        &\Code{split}:
        \\&\quad \ul{l}{:}\rt{il}{\I{len}(\vnu) \geq 4} \sarr
        \\&\quad \ol{\Code{l1}}{:}\rt{il}{\I{len}(\vnu) \geq 2} \times \urt{il}{\I{len}(\vnu) = 2}
    \end{align*}\noindent
    This type can fit both implementations of the $\Code{split}$ function shown in Figure \ref{fig:split}.
    \item Intersection type: The recursive function $\Code{merge}$ also requires correctness specifications (otherwise, recursive calls cannot be typechecked):
    \begin{align*}
        &\Code{merge}:
        \\&\quad \ol{\Code{l1}}{:}\rt{il}{\I{len}(\vnu) < 2 \land \I{sorted}(\vnu)} \sarr \ol{\Code{l2}}{:}\rt{il}{\I{len}(\vnu) < 2 \land \I{sorted}(\vnu)} \sarr 
        \\&\quad \rt{il}{\I{sorted}(\vnu) \land \I{len}(\vnu) = \I{len}(\Code{l1}) + \I{len}(\Code{l2}) \land \I{mem}\ldots}
    \end{align*}\noindent
    Thus, the $\Code{merge}$ function should be annotated with intersected types.
\end{enumerate}

\section{Type Semantics}

In order to analyze reachability, we need both under and overapproximation (e.g., may-must analysis\cite{MayMust}).

\paragraph{Base Type} provide both an underbound and an overbound for the base type.
\begin{align*}
    \denot{\phi} &\doteq \{ v ~|~ \phi(v) \}
    \\ \denot{e} &\doteq \{ v ~|~ \mstep{e}{v} \}
    \\ \denot{\rt{b}{\phi ~|~ \psi}} &\doteq (\denot{\phi}, \denot{\psi})
    \\ e \text{ in denotation of $\tau$ means } &
    (U, O) = \denot{\tau}, U \subseteq \denot{e} \subseteq O
\end{align*}\noindent

\paragraph{Function Type} What happens when a function $e$ apply to a term that has type $\rt{b}{\phi ~|~ \psi}$?
\begin{align*}
    &\text{over application: a type try to apply to $\ort{b}{\psi}$}
    \\ \denot{x{:}\rt{b}{\psi} \sarr \tau} &\doteq \{ e ~|~ \forall v. v \in \denot{\phi} \impl \I{fst}(\denot{\tau[x\mapsto v]}) \subseteq \denot{e\;v} \subseteq \I{snd}(\denot{\tau[x\mapsto v]}) \}
\end{align*}\noindent
Underapproximation? try to describe the lower and higher bound of a term $e\;(\Code{randomPick}\;\denot{\phi})$.
\begin{align*}
    &\text{under application: a type try to apply to $\urt{b}{\phi}$}
    \\ \denot{x{:}\rt{b}{\phi} \sarr \tau} &\doteq \{ e ~|~ \bigcup_{v \in \denot{\phi}} \I{fst}(\denot{\tau[x\mapsto v]}) \subseteq \bigcup_{v \in \denot{\phi}} \denot{e\;v} \subseteq \bigcup_{v \in \denot{\phi}} \I{snd}(\denot{\tau[x\mapsto v]}) \}
\end{align*}\noindent
Simulate the Incorrectness style reverse definition. Intuition, hide ghost varaible (e.g., the actual parameter value) as existential variable (variable in big union).

\paragraph{Typechecking of function}
\begin{align*}
    &\Gamma \vdash \zzlam{x}{e} : x{:}\ort{b}{\phi} \sarr \tau
    \\&\Gamma, x{:}\rt{b}{ \diamond ~|~ \phi} \vdash e : \tau
\end{align*}\noindent
The $\diamond$ means ``exists a value $v$, that is reachable (i.e., $\exists v. \rt{b}{ \vnu = v ~|~ \phi}$)''
\begin{align*}
    &\Gamma \vdash \zzlam{x}{e} : x{:}\urt{b}{\phi} \sarr \tau
    \\&\Gamma, x{:}\rt{b}{ \phi ~|~ \phi} \vdash e : \tau
\end{align*}\noindent
Here the upper bound is the same as the lower bound, because we don't case any other input out of the range of $\phi$.

\paragraph{Typechecking of function application}
\begin{align*}
    f{:}a{:}\ort{b}{\phi}\sarr\tau, x{:}\rt{b}{\phi' ~|~ \psi} &\vdash f\;x : ?
    \\f{:}a{:}\ort{b}{\phi}\sarr\tau, x{:}\rt{b}{\phi' \land \phi ~|~ \psi \land \phi} &\vdash f\;x : \tau[a\mapsto x]
    \\\text{Optimization (prune):}& \text{check that both $\phi' \land \phi$ and $\psi \land \phi$ are not empty}
    \\& \diamond \land \phi = \diamond
\end{align*}\noindent
Underapproximation? The hidden ghost variable is ``consumed''.
\begin{align*}
    f{:}a{:}\urt{b}{\phi}\sarr\tau, x{:}\rt{b}{\phi' ~|~ \psi} &\vdash f\;x : ?
    \\f{:}a{:}\urt{b}{\phi}\sarr\tau, x{:}\rt{b}{\diamond ~|~ \phi} &\vdash f\;x : \tau[a\mapsto x]
    \\\text{check: }& \phi \subseteq \phi' \subseteq \psi
\end{align*}\noindent

\paragraph{Typecontext partitioning}
\begin{align*}
\Gamma, x{:}(\tau_1 \unionty \tau_2) &= \Gamma, x{:}\tau_1 \otimes \Gamma, x{:}\tau_2 \\
\Gamma, x{:}\rt{b}{\phi_1 \cup \phi_2 ~|~ \psi} &= \Gamma, x{:}\rt{b}{\phi_1 ~|~ \psi \land \neg \phi_2} \otimes \Gamma, x{:}\rt{b}{\phi_2 ~|~ \psi \land \neg \phi_1}
\end{align*}\noindent


\paragraph{Union Type and Intersection Type}
\begin{align*}
    \denot{\urt{b}{\phi_1}} \cup \denot{\urt{b}{\phi_2}} \text{ is not closed in qualifier logic}
\end{align*}\noindent

\paragraph{Denotation of Types}
\begin{align*}
    \denot{\tau} &\doteq \ol{\denot{\tau}} \cap \ul{\denot{\tau}}
    \\\ol{\denot{\rt{b}{\phi ~|~ \psi}}} &\doteq \{e ~|~ \denot{e} \subseteq \denot{\psi} \}
    \\ \ul{\denot{\rt{b}{\phi ~|~ \psi}}} &\doteq \{e ~|~ \denot{\phi} \subseteq \denot{e} \}
    \\ \denot{\ol{x}{:}\tau_x \sarr \tau} &\doteq \{e ~|~ \forall v. v \in \ol{\denot{\tau_x}} \impl \ul{\denot{\tau[x\mapsto v]}} \subseteq \denot{e\;v} \subseteq \ol{\denot{\tau[x\mapsto v]}} \}
    \\ \denot{\ul{x}{:}\tau_x \sarr \tau} &\doteq \{e ~|~ \bigcup_{v \in \ul{\denot{\tau_x}}} \ul{\denot{\tau[x\mapsto v]}} \subseteq \bigcup_{v \in \ul{\denot{\tau_x}}} \denot{e\;v} \subseteq \bigcup_{v \in \ul{\denot{\tau_x}}} \ol{\denot{\tau[x\mapsto v]}} \}
\end{align*}\noindent


\section{Algorithm: Higher-Order Model Checking}

Consider the program shown in Figure \ref{fig:higher-order-model-checking}, we want to prove the program will eventually clase the file or loop forever.

\begin{figure}[!t]
    \begin{minted}[xleftmargin=5pt, numbersep=4pt, linenos = true, fontsize = \footnotesize, escapeinside=??]{OCaml}
    let rec g (x: bool) : bool =
      if bool_gen () then false else let x = read x; g x
    in g true
    ?$: \rt{\Bool}{\bot | \vnu = \Code{false}}$?
    \end{minted}
    \caption{The higher-order model checking.}
    \label{fig:higher-order-model-checking}
\end{figure}

The rule of recursion scheme is program in cps:
\begin{align*}
    &\Code{Start} \to g\;true
    \\& g\;x \to \Code{ifElseThen}\; \Code{close}(x)\; (g\; \Code{read}(x))
\end{align*}\noindent
The transition rules of automaton are types (semantics) of builtin operators.
\begin{align*}
    & \vdash \Code{ifElseThen} : \ort{\Bool}{\vnu = \Code{false}}\sarr \ort{\Bool}{\vnu = \Code{false}} \sarr \rt{\Bool}{\vnu = \Code{false} | \vnu = \Code{false}}
    \\& \vdash \Code{ifElseThen} : \ort{\Bool}{\vnu = \Code{true}}\sarr \ort{\Bool}{\vnu = \Code{true}} \sarr \rt{\Bool}{\vnu = \Code{true} | \vnu = \Code{true}} \quad \text{overapproximation style}
    \\& \vdash \Code{read} : \ort{\Bool}{\vnu = \Code{true}}\sarr \rt{\Bool}{\vnu = \Code{true} | \vnu = \Code{true}}
    \\& \vdash \Code{close} : \ort{\Bool}{\vnu = \Code{true}}\sarr \rt{\Bool}{\vnu = \Code{false} | \vnu = \Code{false}}
    \\& \quad \text{\Code{read} and \Code{close} are only valid when the argument is \Code{true}}
\end{align*}\noindent
The proof goal is
\begin{align*}
    \vdash \Code{Start} : \rt{\Bool}{\bot | \vnu = \Code{false}}
\end{align*}\noindent
We need to infer that
\begin{align*}
    &\vdash g: \ort{\Bool}{\vnu = \Code{false} \lor \vnu = \Code{true}}\sarr \rt{\Bool}{\bot | \vnu = \Code{false}}
\end{align*}\noindent

\paragraph{Approach} Note that there are only two states:
\begin{align*}
    &q_0 \doteq \rt{\Bool}{\bot | \vnu = \Code{false}} \quad \text{loop forever or ends with false}
    \\&q_1 \doteq : \rt{\Bool}{\bot | \vnu = \Code{true}} \quad \text{loop forever or ends with false}
\end{align*}\noindent
Which are sufficient to describe builtin operators.
\begin{align*}
    & \vdash \Code{ifElseThen} : q_0 \sarr q_0 \sarr q_0
    \\& \vdash \Code{ifElseThen} : q_1 \sarr q_1 \sarr q_1
    \\& \vdash \Code{read} : q_1 \sarr q_1
    \\& \vdash \Code{close} : q_1 \sarr q_0
\end{align*}\noindent
Try to enumerate all possible types for $\Code{Start}$ and $g$.
\begin{align*}
    &\Code{Start} : q_0, : q_1, 
    \\&g : q_0 \sarr q_0, q_1 \sarr q_1, q_0 \sarr q_1, q_1 \sarr q_0, q_0 \land q_1 \sarr q_0, q_0 \land q_1 \sarr q_1, q_0 \lor q_1 \sarr q_0, q_0 \lor q_1 \sarr q_1
\end{align*}\noindent
Note, the algorithm requires more precise intersection types within parameter types (not for return type, maybe because it can be derived by typechecking a term for multiple times):
\begin{align*}
    &q_0 \doteq \rt{\Bool}{\bot | \vnu = \Code{false}} \quad \text{loop forever or ends with false}
    \\&q_1 \doteq : \rt{\Bool}{\bot | \vnu = \Code{true}} \quad \text{loop forever or ends with false}
    \\&q_0 \land q_1 \doteq : \rt{\Bool}{\bot | \bot} \quad \text{loop forever}
    \\&q_0 \lor q_1 : \rt{\Bool}{\bot | \top} \quad \text{any possible state}
\end{align*}\noindent
Then we calculate a greatest fixpoint of these types for $\Code{Start}$ and $g$. For example, assume a type context $\Gamma$ contains all possible types for $\Code{Start}$ and $g$, we type check $\Gamma \vdash g: q_1 \sarr q_1$, that is
\begin{align*}
    &\Gamma, x{:}\rt{\Bool}{\vnu = \Code{true} ~|~ \vnu = \Code{true}} \vdash g: \rt{\Bool}{\bot ~|~ \vnu = \Code{true}}
    \\&\quad \text{need to prove the following two judgements:}
    \\&\Gamma, x{:}\rt{\Bool}{\vnu = \Code{true} ~|~ \vnu = \Code{true}} \vdash \Code{false} : \rt{\Bool}{\bot ~|~ \vnu = \Code{true}}
    \\&\Gamma, x{:}\rt{\Bool}{\vnu = \Code{true} ~|~ \vnu = \Code{true}} \vdash g\; \Code{read}(x) : \rt{\Bool}{\bot ~|~ \vnu = \Code{true}}
    \\&\quad \text{The first fails, thus we remove $g: q_1 \sarr q_1$ from $\Gamma$. }
\end{align*}\noindent

\section{Algorithm: Higher-Order Symbolic Execution}

\begin{figure}[!t]
    \begin{minted}[xleftmargin=5pt, numbersep=4pt, linenos = true, fontsize = \footnotesize, escapeinside=??]{OCaml}
    let f (g: int -> int) (n: int) : int =
      1 / (100 - (g n))
    in ? f
    \end{minted}
    \caption{The higher-order symbolic execution.}
    \label{fig:higher-order-symbolic-execution}
\end{figure}

\paragraph{Example} Consider the program shown in Figure \ref{fig:higher-order-symbolic-execution}, we want to find a valid instantation of function $?$ such that leads division by zero exception. It is equivalent to find a valid type of function $h$ such that the program type checks.
\begin{align*}
    &\vdash (/) : \Int \sarr \ort{\Int}{\vnu \neq 0} \sarr \Int
    \\&\vdash (-) : x{:}\Int \sarr y{:}\ort\Int \sarr \rt{\Int}{\vnu = x - y ~|~ \vnu = x - y}
    \\&h{:}((\Int\sarr\Int)\sarr\Int\sarr\Int)\sarr\Int \vdash h\; \zzlam{g}{\zzlam{n}{1 /_{\Code{Err}} (100 - (g\; n))}} : \Code{Err}_{\Int}
\end{align*}\noindent
The paper considers $3$ situations for the function application $h\;f$ (when $f$ is a function):
\begin{enumerate}
    \item the body of $h$ has no reference to $f$.
    \item the body of $h$ provides a concrete argument of $f$,
    \item the body of $h$ doesn't provide a concrete argument of $f$, e.g., $\zzlam{f}{f}$.
\end{enumerate}
For the first two situations, the application cannot raise a division by zero exception at location $/_{\Code{Err}}$. For the third situation, it doesn't consistent with the type of $h$.

\paragraph{Approach} We disallowing $\zlet{x}{v}{e}$, thus we see a application $h\;f$ where $h$ is a varaible iff we don't know the implementation of $h$.
\begin{figure}[!t]
{\small
\mprooftr{50mm}
  {}
  {AppFreeVar}
  {$\Gamma, h{:}\tau_h \vdash h\;v : \tau \typeinfer \Gamma, h{:}\tau_h \interty (t_g \sarr \tau)$}
\\[.8em]
\mprooftr{70mm}
  {$\Gamma, h{:}\tau_h \interty ((g{:}\urt{b}{\top} \sarr \Code{Err}) \sarr \tau), g{:}\urt{b}{\top} \vdash e : \Code{Err}$ \quad $\Code{Err} \in \tau$}
  {AppAppErr}
  {$\Gamma, h{:}\tau_h \vdash h\;\zzlam{g}{e} : \tau \typeinfer \Gamma, h{:}\tau_h \interty ((g{:}\urt{b}{\top} \sarr \Code{Err}) \sarr \tau)$}
}
\end{figure}

For example,
\begin{align*}
    &\vdash h\; \zzlam{g}{\zzlam{n}{1 /_{\Code{Err}} (100 - (g\; n))}} : \Code{Err}_{\Int} \tag{$\textsc{AppAppErr}$}
    \\&h{:}..., \ul{g}{:}(y{:}\ort{\Int}{\top}\sarr\rt{\Int}{\bot ~|~ \top}), \ul{n}{:}\rt{\Int}{\bot ~|~ \top} \vdash 1 /_{\Code{Err}} (100 - (g\; n)) : \Code{Err}_{\Int} \tag{$\textsc{AppAppErr}$}
    \\&h{:}..., \ul{g}{:}(y{:}\ort{\Int}{\top}\sarr\rt{\Int}{\bot ~|~ \top}), \ul{n}{:}\rt{\Int}{\bot ~|~ \top} \vdash 100 - (g\; n) : \rt{\Int}{\vnu = 0 ~|~ \vnu = 0}
    \\&h{:}..., \ul{g}{:}(y{:}\ort{\Int}{\top}\sarr\rt{\Int}{\bot ~|~ \top}), \ul{n}{:}\rt{\Int}{\bot ~|~ \top} \vdash g\; n : \rt{\Int}{\vnu = 100 ~|~ \vnu = 100}
    \\&h{:}..., \ul{g}{:}(y{:}\ort{\Int}{\top}\sarr\rt{\Int}{\vnu = 100 ~|~ \vnu = 100}), \ul{n}{:}\rt{\Int}{\bot ~|~ \top} \vdash g\; n : \rt{\Int}{\vnu = 100 ~|~ \vnu = 100} \tag{$\textsc{AppFreeVar}$}
\end{align*}\noindent
